\expandafter\ifx\csname CAPETINCLUDEMODE\endcsname\relax
\documentclass[11pt,a4paper]{article}
% preamble.tex — packages et commandes communes a toutes les fiches CAPET SII IE
\input{glyphtounicode}
\pdfgentounicode=1
\usepackage{cmap}
\usepackage[T1]{fontenc}
\usepackage[utf8]{inputenc}
\usepackage[french]{babel}
% babel[french] active a \begin{document} des raccourcis de ponctuation (:;!?)
% qui cassent la syntaxe to[...] de circuitikz (erreur "+ or - expected") ;
% on les desactive, ce qui ne change que l'espacement fin avant ces signes.
\AtBeginDocument{\shorthandoff{:;!?}}
\usepackage[a4paper, top=1.6cm, bottom=1.6cm, left=1.4cm, right=1.4cm]{geometry}
\setlength{\headheight}{14.5pt}
\usepackage{amsmath, amssymb, amsfonts}
\usepackage{siunitx}
\sisetup{locale=FR, output-decimal-marker={,}, inter-unit-product=\ensuremath{{}\cdot{}},
         group-minimum-digits=4}
% Unites non fournies par siunitx, declarees ici une fois pour toutes (fiches + exercices).
% Les gardes \ifdefined evitent tout conflit avec une redeclaration locale restee dans un
% fichier : ne pas retirer les gardes, ne pas redeclarer ces unites dans les fichiers.
\DeclareSIUnit{\jour}{j}
\ifdefined\var\else\DeclareSIUnit{\var}{var}\fi   % puissance reactive
\ifdefined\vad\else\DeclareSIUnit{\vad}{VAD}\fi   % puissance deformante
\ifdefined\tr\else\DeclareSIUnit{\tr}{tr}\fi      % tour (tr/min)
\usepackage{array, booktabs, multirow, tabularx}
\usepackage{xcolor}
\usepackage[most, breakable]{tcolorbox}
\usepackage{multicol}
\usepackage{enumitem}
\usepackage{mathtools}
\usepackage{lmodern}
% Refonte visuelle 2026 (direction "moderne coloré") : Lato pour le texte,
% maths sans-serif assorties via newtxsf (charge apres amsmath/amssymb).
\usepackage[default]{lato}
\usepackage{newtxsf}
\usepackage{microtype}
\usepackage{fancyhdr}
\usepackage{tikz}
\usetikzlibrary{positioning,arrows.meta}
\usepackage{circuitikz}
\usepackage{lastpage}

% ── Couleurs par domaine (palette modernisée 2026) ────────────────────
\definecolor{colElec}{RGB}{29, 78, 216}        % Électrotechnique — bleu vif
\definecolor{colElecLight}{RGB}{219, 234, 254} % fond bleu clair
\definecolor{colAuto}{RGB}{21, 128, 61}        % Automatique — vert
\definecolor{colMeca}{RGB}{194, 65, 12}        % Mécanique — orange
\definecolor{colTherm}{RGB}{185, 28, 28}       % Thermique — rouge
\definecolor{colQT}{RGB}{232, 244, 233}        % (hérité, plus utilisé par questiontype)
\definecolor{colJury}{RGB}{255, 248, 225}      % fond jury (crème)
\definecolor{colJuryFrame}{RGB}{202, 138, 4}   % bordure/bandeau jury (ambre)
\definecolor{colGrayLight}{RGB}{245, 245, 247} % fond gris neutre
\definecolor{colSeparator}{RGB}{180, 180, 180} % gris séparateur
\definecolor{colExo}{RGB}{8, 145, 178}          % Banque d'exercices — cyan (distinct des fiches)
% ── Encre et gris de la refonte ───────────────────────────────────────
\definecolor{colInk}{RGB}{23, 32, 42}          % texte courant (quasi-noir)
\definecolor{colSlate}{RGB}{83, 98, 118}       % gris ardoise (méta, en-têtes)
\definecolor{colRule}{RGB}{203, 213, 225}      % filets clairs
\definecolor{colJuryBg}{RGB}{255, 248, 225}    % fond crème jury
\definecolor{colPresence}{RGB}{83, 98, 118}    % bloc 5 (présence dans les sujets)
\AtBeginDocument{\color{colInk}}

% ── Style encadré exercice + corrigé (banque d'exercices) ─────────────
\tcbset{
  exoBox/.style={
    breakable,
    enhanced,
    colback=colExo!8!white,
    colframe=colExo,
    colbacktitle=colExo,
    coltitle=white,
    fonttitle=\bfseries\small,
    left=4pt, right=4pt, top=3pt, bottom=3pt,
    before skip=5pt, after skip=5pt,
  }
}

% ── Bascule enonce / corrige (banque d'exercices) ─────────────────────
% \CAPETCORRIGE est defini en ligne de commande (pdflatex "\def\CAPETCORRIGE{1}\input{...}")
% pour compiler la variante corriges d'un meme fichier source ; par defaut (non defini),
% la compilation produit la variante enonces. \CAPETSHOW{enonce}{corrige} choisit la branche.
%
% TROIS MODES (refonte 2026-09-20, manuel de revision) :
%   - aucun drapeau   -> imprime l'enonce            (variante enonces, INCHANGE)
%   - \CAPETCORRIGE   -> imprime le corrige          (variante corriges, INCHANGE)
%   - \CAPETMANUEL    -> imprime l'enonce ET met le corrige de cote pour l'annexe
%
% Le mode manuel stocke chaque corrige dans une macro (\csgdef), sans ecriture
% fichier ni re-tokenisation : le contenu (circuitikz, accents, %) traverse intact.
% L'annexe les rejoue ensuite dans l'ordre via \CAPETANNEXE.
\usepackage{etoolbox}

\newcounter{CAPETexo}          % numero continu d'exercice dans le manuel
\setcounter{CAPETexo}{0}
\newcounter{CAPETannexe}       % compteur de relecture en annexe

% Titre courant de la notion en cours (pose par \CAPETNOTION), sert d'intertitre annexe
\newcommand{\CAPETnotioncode}{}
\newcommand{\CAPETnotiontitre}{}

% \CAPETNOTION{code}{titre} : ouvre une notion dans le manuel (memorise le libelle,
% et marque le rang du premier exercice de la notion pour l'intertitre d'annexe).
\newcommand{\CAPETNOTION}[2]{%
  \renewcommand{\CAPETnotioncode}{#1}%
  \renewcommand{\CAPETnotiontitre}{#2}%
  \csgdef{CAPETnotionstart\the\numexpr\value{CAPETexo}+1\relax}{#1 -- #2}%
}

% ── \degree robuste ───────────────────────────────────────────────────
% Certains fichiers d'exercices utilisent \degree. Ni le preambule ni siunitx
% ne le definissent, et \providecommand local ne suffit pas : sous hyperref,
% \degree existe deja comme commande factice (\ERROR), donc \providecommand
% s'abstient et l'appel echoue a la compilation du manuel (erreur invisible
% en compilation isolee). On le (re)definit donc inconditionnellement ici,
% apres le chargement des packages, pour tout le corpus.
\AtBeginDocument{\providecommand{\degree}{}\renewcommand{\degree}{\ensuremath{^\circ}}}

% Libelle decoratif : les fichiers d'exercices utilisent aussi \CAPETSHOW dans
% leur \fancyhead et leur bloc de titre, pour afficher "(enonces)" ou "(corriges)".
% Ces appels-la ne sont PAS des exercices : en mode manuel ils ne doivent ni
% consommer un numero, ni alimenter l'annexe. On les reconnait a leur premier
% argument, strictement egal au libelle ci-dessous.
\newcommand{\CAPETLIBELLEENONCES}{(\'enonc\'es)}

\ifdefined\CAPETMANUEL
  % --- Mode manuel : enonce en place, corrige differe vers l'annexe ---
  \newcommand{\CAPETSHOW}[2]{%
    \def\CAPETargun{#1}%
    \ifx\CAPETargun\CAPETLIBELLEENONCES
      #1%            appel decoratif : rendu tel quel, hors numerotation
    \else
      \CAPETSHOWEXO{#1}{#2}%
    \fi
  }
  \newcommand{\CAPETSHOWEXO}[2]{%
    \stepcounter{CAPETexo}%
    \csgdef{CAPETcorrige\theCAPETexo}{#2}%
    \edef\CAPETtmpnum{\theCAPETexo}%
    #1%
    \par\vspace{2pt}%
    {\footnotesize\itshape\color{colExo}%
     \hyperref[capet:cor:\CAPETtmpnum]{Corrige de l'exercice \CAPETtmpnum\ $\rightarrow$ p.~\pageref{capet:cor:\CAPETtmpnum}}%
     \label{capet:eno:\CAPETtmpnum}}%
  }
  % \CAPETANNEXE : rejoue tous les corriges captures, dans l'ordre, avec intertitres
  \newcommand{\CAPETANNEXE}{%
    \setcounter{CAPETannexe}{0}%
    \loop\ifnum\value{CAPETannexe}<\value{CAPETexo}%
      \stepcounter{CAPETannexe}%
      \edef\CAPETcurnum{\theCAPETannexe}%
      \ifcsname CAPETnotionstart\CAPETcurnum\endcsname
        \par\medskip
        {\normalsize\bfseries\color{colExo}\csuse{CAPETnotionstart\CAPETcurnum}}%
        \par\nobreak\smallskip
      \fi
      \begin{tcolorbox}[exoBox, breakable,
        title={Corrige de l'exercice \CAPETcurnum\
               \hfill\normalfont\footnotesize
               \hyperref[capet:eno:\CAPETcurnum]{enonce p.~\pageref{capet:eno:\CAPETcurnum}}}]%
      \label{capet:cor:\CAPETcurnum}%
      \csuse{CAPETcorrige\CAPETcurnum}%
      \end{tcolorbox}%
    \repeat
  }
\else
  % --- Modes historiques, comportement strictement inchange ---
  \newcommand{\CAPETSHOW}[2]{\ifdefined\CAPETCORRIGE #2\else #1\fi}
  \newcommand{\CAPETANNEXE}{}
\fi

% ── Difficulté 1 à 3 (étoiles pleines, amssymb) ───────────────────────
\newcommand{\difficulte}[1]{%
  \ifcase#1\relax
  \or $\bigstar$\;niveau 1/3%
  \or $\bigstar\bigstar$\;niveau 2/3%
  \or $\bigstar\bigstar\bigstar$\;niveau 3/3%
  \fi
}

% ── Style encadré extrait de document technique (banque d'exercices) ──
\tcbset{
  docTech/.style={
    breakable,
    enhanced,
    colback=colGrayLight,
    colframe=colSeparator,
    fonttitle=\bfseries\small\itshape,
    coltitle=black,
    colbacktitle=colGrayLight,
    left=4pt, right=4pt, top=3pt, bottom=3pt,
    before skip=4pt, after skip=4pt,
  }
}

% ── Style tcolorbox de section principale ────────────────────────────
% Refonte 2026 : bandeau titre coloré arrondi + fond très pâle, sans cadre.
\tcbset{
  ficheSection/.style={
    breakable, enhanced jigsaw,
    frame hidden, boxrule=0pt, arc=1.8mm,
    colback=#1!4!white,
    colbacktitle=#1,
    coltitle=white,
    fonttitle=\sffamily\bfseries,
    left=7pt, right=7pt, top=5pt, bottom=5pt,
    toptitle=3pt, bottomtitle=3pt,
    before skip=10pt, after skip=8pt,
  }
}

% ── Style encadré question type ──────────────────────────────────────
% Carte blanche, bandeau sombre, liseré très clair.
\tcbset{
  questiontype/.style={
    breakable, enhanced jigsaw,
    colback=white,
    frame hidden, boxrule=0pt, arc=1.8mm,
    colbacktitle=colElec!78!black,
    coltitle=white,
    fonttitle=\sffamily\bfseries\small,
    left=7pt, right=7pt, top=4pt, bottom=4pt,
    toptitle=2.5pt, bottomtitle=2.5pt,
    before skip=8pt, after skip=8pt,
    borderline={0.6pt}{0pt}{colElec!12!white},
  }
}

% ── Style encadré jury ───────────────────────────────────────────────
% Bandeau ambre, fond crème, sans cadre.
\tcbset{
  jury/.style={
    breakable, enhanced jigsaw,
    colback=colJuryBg,
    frame hidden, boxrule=0pt, arc=1.8mm,
    colbacktitle=colJuryFrame,
    coltitle=white,
    fonttitle=\sffamily\bfseries,
    left=7pt, right=7pt, top=4pt, bottom=4pt,
    toptitle=3pt, bottomtitle=3pt,
    before skip=10pt, after skip=8pt,
  }
}

% ── En-tête / pied de page communs (refonte 2026) ─────────────────────
% #1 code fiche, #2 titre court, #3 couleur domaine
\newcommand{\ficheHeaders}[3]{%
  \fancyhead[L]{\small\sffamily{\bfseries\textcolor{#3}{#1}}\hspace{0.7em}\textcolor{colSlate}{#2}}%
  \fancyhead[R]{\small\sffamily\textcolor{colSlate}{CAPET SII IE\, --- \,Session 2026}}%
  \fancyfoot[C]{\small\sffamily\textcolor{colSlate}{\thepage\,/\,\pageref{LastPage}}}%
  \fancyhead[C]{}\fancyfoot[L]{}\fancyfoot[R]{}%
}

% ── Bloc en-tête de fiche (BLOC 0, refonte 2026) ──────────────────────
% #1 code  #2 titre  #3 couleur domaine  #4 domaine  #5 ligne présence  #6 sous-thèmes
\newcommand{\ficheEntete}[6]{%
  \begin{tcolorbox}[enhanced, breakable,
    frame hidden, boxrule=0pt, arc=2.5mm,
    interior style={left color=#3, right color=#3!75!black},
    left=9pt, right=9pt, top=7pt, bottom=7pt,
    before skip=0pt, after skip=6pt]
    {\color{white}\sffamily\bfseries\LARGE #1\quad #2}\\[4pt]
    {\color{#3!22!white}\sffamily\small \MakeUppercase{#4}\quad\textbar\quad #5}\\[1pt]
    {\color{#3!35!white}\sffamily\footnotesize Sous-thèmes : #6}
  \end{tcolorbox}%
}

% ── Mise en page ─────────────────────────────────────────────────────
\pagestyle{fancy}
\fancyhf{}
\renewcommand{\headrulewidth}{0.5pt}
\renewcommand{\footrulewidth}{0pt}
\AtBeginDocument{%
  \renewcommand{\headrule}{{\color{colRule}\hrule width\headwidth height 0.5pt}}%
}

\setlist[itemize]{leftmargin=1.2em, itemsep=1pt, topsep=2pt, parsep=0pt}
\setlist[enumerate]{leftmargin=1.2em, itemsep=3pt, topsep=2pt, parsep=0pt}

\setlength{\parskip}{2pt}
\setlength{\columnsep}{10pt}

\def\CAPETBEGINFILE{\begin{document}}
\def\CAPETENDFILE{\end{document}}
\else
\def\CAPETBEGINFILE{}
\def\CAPETENDFILE{}
\fi

\ficheHeaders{T08}{Redresseurs}{colElec}

\CAPETBEGINFILE

% -----------------------------------------------------------------------
%  BLOC EN-TETE
% -----------------------------------------------------------------------

\ficheEntete{T08}{Redresseurs (monophasé et triphasé)}{colElec}%
  {électrotechnique - électronique de puissance}%
  {\textbf{Présence :} \textbf{13/13} sujets (2013-2026) - thème \emph{omniprésent}}%
  {PD2 (Graëtz monophasé) - P3 (triphasé simple alt.) - PD6 (triphasé double alt.) - facteur de puissance - harmoniques - redresseur commandé (thyristors)}

% -----------------------------------------------------------------------
%  1. RAPPELS DE COURS
% -----------------------------------------------------------------------

\begin{tcolorbox}[
  ficheSection=colElec,
  title={1. Rappels de cours},
]

% -- 1.1 Tensions moyennes ---------------------------------------------
\textbf{1.1 Tensions moyennes de sortie}

\smallskip
\begin{center}
\renewcommand{\arraystretch}{1.5}
\begin{tabular}{lllll}
\toprule
\textbf{Pont} & \textbf{Diodes} & \textbf{Formule} & \textbf{Valeur num.} & \textbf{Pulsations/T} \\
\midrule
PD2 (monophasé) & 4 & $\bar{U} = \dfrac{2\sqrt{2}}{\pi} V_{eff}$ & $\approx 0{,}9\,V_{eff}$ & 2 \\[8pt]
P3 (triphasé) & 3 & $\bar{U} = \dfrac{3\sqrt{6}}{2\pi} V_{eff,simple}$ & $\approx 1{,}17\,V_{s,simple}$ & 3 \\[8pt]
PD6 (triphasé) & 6 & $\bar{U} = \dfrac{3\sqrt{2}}{\pi} V_{eff,comp}$ & $\approx 2{,}34\,V_{s,simple}$ & 6 \\[8pt]
PD6 compact & 6 & $\bar{U} = 0{,}95 \cdot U_M$ & $U_M$ = ampli. composées & 6 \\
\bottomrule
\end{tabular}
\end{center}

\smallskip
\textbf{Chronogramme de la tension redressée (PD2, double alternance) :}
\begin{center}
\begin{tikzpicture}[scale=0.85, font=\footnotesize]
  \draw[->] (0,0) -- (6.5,0) node[right] {$t$};
  \draw[->] (0,-0.3) -- (0,2.3) node[above] {$u_s(t)$};
  \draw[thick, colElec, domain=0:6.28, samples=100, smooth] plot (\x, {1.8*abs(sin(\x r))});
  \draw[dashed] (0,1.15) -- (6.4,1.15) node[right] {\scriptsize $\bar{U}$};
\end{tikzpicture}
\end{center}
Chaque « bosse » correspond à une pulsation : $q=2$ pour PD2 (visible ci-dessus), $q=6$ pour PD6 - un PD6 lisse davantage la tension et réduit le taux d'ondulation.

\medskip
\hrule\vspace{4pt}

% -- 1.2 Facteur de puissance -----------------------------------------
\textbf{1.2 Facteur de puissance}

\smallskip
Pour un redresseur à diodes débitant sur une charge à courant parfaitement lissé, le courant de
ligne n'est pas sinusoïdal ($f_P \neq \cos\varphi$). Les diodes n'étant pas commandées, le
fondamental du courant est en phase avec la tension ($\cos\varphi_1 = 1$) et le facteur de
puissance se réduit au \textbf{facteur de distorsion} :
\[
  f_P = \frac{I_{1,\text{eff}}}{I_{\text{eff}}}\,\cos\varphi_1 = \frac{I_{1,\text{eff}}}{I_{\text{eff}}}
  \qquad\text{ou, de façon équivalente,}\qquad f_P = \frac{P_A}{S} = \frac{\langle u_S\rangle\,I_S}{S}
\]
\textbf{Il n'existe pas de formule unique en $q$ valable pour tous les montages} : chaque topologie
impose sa propre durée de conduction, donc son propre $I_{\text{eff}}$. Voir le tableau ci-dessous
pour les trois cas usuels, et la Question Type 1 pour l'établissement complet dans le cas du P3.

\begin{center}
\renewcommand{\arraystretch}{1.35}
\begin{tabular}{llll}
\toprule
\textbf{Pont} & $q$ & $f_P$ & \textbf{Conformité CDC $f_P \geq 0{,}9$} \\
\midrule
PD2 & 2 & $\dfrac{2\sqrt{2}}{\pi} \approx 0{,}90$ & Limite \\[6pt]
P3  & 3 & $\dfrac{3\sqrt{2}}{2\pi} \approx 0{,}675$ & Non \\[6pt]
PD6 & 6 & $\dfrac{3}{\pi} \approx 0{,}955$ & \textbf{Oui} \\
\bottomrule
\end{tabular}
\end{center}

\medskip
\hrule\vspace{4pt}

% -- 1.3 Harmoniques injectes -----------------------------------------
\textbf{1.3 Harmoniques côté réseau}

\smallskip
Les redresseurs à diodes injectent des harmoniques de courant côté réseau.
Rangs présents : $6k \pm 1$ pour PD6 ($k = 1, 2, \ldots$) - rangs 5, 7, 11, 13, \ldots
Pour PD2 : rangs $2k \pm 1$ - rangs 3, 5, 7, \ldots\\
\textit{Taux d'ondulation tension : PD6 = 4,2\,\% - P3 = 18,3\,\% - PD2 = 48,2\,\%}

\end{tcolorbox}

% -----------------------------------------------------------------------
%  2. FORMULAIRE SYNTHETIQUE
% -----------------------------------------------------------------------

\begin{tcolorbox}[
  ficheSection=colElec,
  title={2. Formulaire synthétique},
]
\small
\begin{multicols}{2}
\renewcommand{\arraystretch}{1.45}
\begin{tabular}{@{}p{4.6cm}ll@{}}
\toprule
\textbf{Formule} & \textbf{Usage} & \textbf{U.} \\
\midrule
\multicolumn{3}{@{}l}{\footnotesize\textit{Tensions moyennes}} \\
$\bar{U}_{PD2} = 0{,}9\,V_{eff}$ & Monoph. & V \\
$\bar{U}_{P3} = 1{,}17\,V_{s,simple}$ & 3 diodes & V \\
$\bar{U}_{PD6} = 2{,}34\,V_{s,simple}$ & 6 diodes & V \\
$\bar{U}_{PD6} = 0{,}95\,U_M$ & Ampli comp. & V \\
$U_{s,comp} = U_M/\sqrt{2}$ & Valeur eff. & V \\
\midrule
\multicolumn{3}{@{}l}{\footnotesize\textit{Rapport de transfo (amont)}} \\
$m = U_{s,comp}/U_{1,comp}$ & Yy & - \\
$m = V_{s,simple}/V_{1,simple}$ & Yy équiv. & - \\
\bottomrule
\end{tabular}

\columnbreak

\renewcommand{\arraystretch}{1.45}
\begin{tabular}{@{}p{4.6cm}ll@{}}
\toprule
\textbf{Formule} & \textbf{Usage} & \textbf{U.} \\
\midrule
\multicolumn{3}{@{}l}{\footnotesize\textit{Facteur de puissance}} \\
$f_P = \frac{I_{1,\text{eff}}}{I_{\text{eff}}}\cos\varphi_1$ & Général & - \\
$f_{P,PD2} = \frac{2\sqrt{2}}{\pi} \approx 0{,}90$ & PD2 & - \\
$f_{P,P3} = \frac{3\sqrt{2}}{2\pi} \approx 0{,}675$ & P3 & - \\
$f_{P,PD6} = \frac{3}{\pi} \approx 0{,}955$ & PD6 & - \\
\midrule
\multicolumn{3}{@{}l}{\footnotesize\textit{Thyristor commandé}} \\
$\bar{U}_\alpha = \bar{U}_0\cos\alpha$ & Pont commandé & V \\
\midrule
\multicolumn{3}{@{}l}{\footnotesize\textit{Courant efficace diode P3}} \\
$I_{D,eff} = I_S/\sqrt{3}$ & Chaque diode & A \\
\bottomrule
\end{tabular}
\end{multicols}
\normalsize
\end{tcolorbox}

% -----------------------------------------------------------------------
%  3. QUESTIONS TYPES
% -----------------------------------------------------------------------

\begin{tcolorbox}[
  ficheSection=colElec,
  title={3. Questions types - extraites des sujets réels},
]

% -- QT1 -------------------------------------------------------------

\begin{tcolorbox}[questiontype,
  title={QT1 - Facteur de puissance P3 vs PD6 - choix de structure \hfill\normalfont\footnotesize E1 sujet 2026}
]
\textbf{Énoncé.}
Un réseau triphasé \SI{230}{\volt}/\SI{400}{\volt} alimente un chargeur de batterie DC. Le cahier des charges impose une tension DC de \SI{600}{\volt} et un facteur de puissance $f_P \geq 0{,}9$. Calculer $f_P$ pour un pont P3 ($q = 3$) et pour un pont PD6 ($q = 6$). Conclure sur la structure à retenir.

\tcbline

\textbf{Corrigé.}

\textbf{Étape 1 — Pont P3.} Chaque diode conduit $I_S$ pendant $T/3$, donc
$I_{\text{eff}} = I_S/\sqrt3$ et $S = 3\,V I_{\text{eff}} = \sqrt3\,V I_S$, tandis que
$P_A = \langle u_S\rangle I_S = \frac{3\sqrt6}{2\pi} V I_S$ :
\[
  f_{P,P3} = \frac{P_A}{S} = \frac{3\sqrt6/(2\pi)}{\sqrt3}
  = \frac{3\sqrt{2}}{2\pi} \approx 0{,}675
\]
$f_{P,P3} = 0{,}675 < 0{,}9$ - \textbf{non conforme}.

\textbf{Étape 2 — Pont PD6.} Chaque enroulement conduit désormais $2T/3$ de façon alternée :
\[
  f_{P,PD6} = \frac{3}{\pi} \approx 0{,}955
\]
$f_{P,PD6} = 0{,}955 > 0{,}9$ - \textbf{conforme}.

\textbf{Conclusion.} Seul le pont PD6 respecte le facteur de puissance minimal imposé ($0{,}955 \geq 0{,}9$) : c'est la structure à retenir pour ce cahier des charges, malgré son nombre de diodes plus élevé.
\end{tcolorbox}

\vspace{4pt}

% -- QT2 -------------------------------------------------------------

\begin{tcolorbox}[questiontype,
  title={QT2 - PD6 à 600\,V DC - calcul du rapport de transformation \hfill\normalfont\footnotesize E1 sujet 2026 (Q.E10)}
]
\textbf{Énoncé.}
Pour obtenir $\bar{u}_s = 600\,\mathrm{V}$ avec un pont PD6 ($\bar{u}_s = 0{,}95 \cdot U_M$) alimenté par un transformateur étoile-étoile (primaire $U_{1,comp} = 400\,\mathrm{V}$), calculer $U_M$, la valeur efficace des tensions composées secondaires $U_{s,comp}$, et le rapport $m$.

\tcbline

\textbf{Corrigé.}

\textbf{Étape 1.}
$U_M = 600/0{,}95 \approx 631\,\mathrm{V}$

\textbf{Étape 2.}
$U_{s,comp} = U_M/\sqrt{2} = 631/1{,}414 \approx 446\,\mathrm{V}$

\textbf{Étape 3.}
\[
  m = \frac{U_{s,comp}}{U_{1,comp}} = \frac{446}{400} \approx \mathbf{1{,}12}
\]

\textbf{Conclusion.} Un rapport de transformation $m \approx 1{,}12$, proche de l'unité, suffit pour atteindre les $600\,\mathrm{V}$ DC visés avec un pont PD6 : le transformateur reste peu dimensionné, ce qui est un avantage économique de cette structure.
\end{tcolorbox}

\vspace{4pt}

% -- QT3 -------------------------------------------------------------

\begin{tcolorbox}[questiontype,
  title={QT3 - PD2 + hacheur - alimentation Urbanloop 230\,V vers 72\,V \hfill\normalfont\footnotesize E2 appliquée 2024}
]
\textbf{Énoncé.}
Le boîtier alimentation de la capsule Urbanloop convertit le réseau monophasé $230\,\mathrm{V}$ en $72\,\mathrm{V}$ DC sur les rails. La structure retenue est : PD2 + filtre LC + hacheur Buck. Calculer $\bar{U}_{PD2}$ puis le rapport cyclique $\alpha$ du hacheur.

\tcbline

\textbf{Corrigé.}

\textbf{Étape 1 — Tension DC du pont PD2.}
\[
  \bar{U}_{PD2} = \frac{2\sqrt{2}}{\pi} \times 230 = 0{,}9 \times 230 \approx 207\,\mathrm{V}
\]

\textbf{Étape 2 — Rapport cyclique du hacheur Buck.}
\[
  \alpha = \frac{\langle u_{sortie} \rangle}{E} = \frac{72}{207} \approx \mathbf{0{,}35}
\]

\textbf{Conclusion.} Le PD2 fixe la tension DC d'entrée ($207\,\mathrm{V}$) et le hacheur Buck ($\alpha \approx 0{,}35$) l'adapte au niveau des rails ($72\,\mathrm{V}$) : cette architecture à deux étages répond au cahier des charges sans recourir à un second transformateur.
\end{tcolorbox}

\end{tcolorbox}

% -----------------------------------------------------------------------
%  4. ATTENTES DU JURY
% -----------------------------------------------------------------------

\begin{tcolorbox}[
  jury,
  title={4. Attentes du jury - Erreurs récurrentes et points de vigilance},
]

\begin{enumerate}[label=\textbf{\arabic*.}, leftmargin=*, itemsep=5pt]

  \item \textbf{Mémoriser les trois formules de $\bar{U}$.}
  PD2 $\approx 0{,}9\,V_{eff}$, P3 $\approx 1{,}17\,V_{s,simple}$, PD6 $\approx 2{,}34\,V_{s,simple}$.
  Les confondre est une erreur éliminatoire dans les questions de dimensionnement.

  \item \textbf{P3 $\neq$ PD6.}
  P3 = 3 diodes (simple alternance, 3 pulsations). PD6 = 6 diodes (pont de Graëtz triphasé, 6 pulsations). Le terme « PD3 » peut apparaître dans certains sujets pour désigner le PD6.

  \item \textbf{Tension PD6 : utiliser les tensions \emph{composées}.}
  $\bar{U}_{PD6} = 0{,}95\,U_M$ avec $U_M$ = amplitude des tensions \textbf{composées}. Utiliser $V_{simple}$ donne une valeur sous-estimée d'un facteur $\sqrt{3}$.

  \item \textbf{Facteur de puissance $\neq$ $\cos\varphi$.}
  Pour un redresseur à diodes, le courant n'est pas sinusoïdal - $f_P \neq \cos\varphi$.
  L'établir par le bilan de puissances ($P_A = \langle u_S\rangle I_S$, $S = 3\,V I_{\text{eff}}$) ou par le facteur de distorsion $I_{1,\text{eff}}/I_{\text{eff}}$ ; se méfier de toute formule unique en $q$, qui ne vaut pas pour tous les montages. Le sujet 2026 ne fournit d'expression directe ($\frac{2}{\pi}\sqrt{q}\sin(\frac{\pi}{q})$, $q$ = nombre de \emph{phases}) que pour le PD3 (Q.E9) ; pour le P3 (Q.E8), elle doit être établie.

  \item \textbf{Conclure par rapport au cahier des charges.}
  Ne pas juste calculer : comparer $f_P$ à la valeur imposée et conclure (conforme / non conforme).

  \item \textbf{Identifier d'abord la tension DC voulue, puis calculer $m$.}
  Le rapport de transformation du transformateur amont est déduit de la tension secondaire AC nécessaire pour obtenir la tension DC cible, pas l'inverse.

\end{enumerate}

\medskip
\begin{tcolorbox}[
  colback=colJuryFrame!12!colJury,
  colframe=colJuryFrame,
  left=4pt, right=4pt, top=2pt, bottom=2pt,
  fontupper=\small\itshape
]
Rapport jury 2025 : «~De nombreux candidats ne maitrisent pas les lois basiques de l'électricité (la loi des nœuds, la loi des mailles, les calculs de puissance ou encore la loi d'Ohm).~» Dans le bilan d'un redresseur, distinguer clairement $P$ (active), $Q$ (réactive), $S$ (apparente) et $f_P = P/S$.\\
Le corrigé du rapport jury 2024 (questions B11-B12, variateur de vitesse) rappelle que la norme IEC/EN 61000-2-2 fixe un THD de courant inférieur à 25\,\% : le rôle du redresseur d'entrée dans le THD global doit être identifié (harmoniques de rangs $2k\pm1$ pour un PD2 côté réseau).
\end{tcolorbox}

\end{tcolorbox}

% -----------------------------------------------------------------------
%  PIED - PRESENCE DANS LES SUJETS
% -----------------------------------------------------------------------

\begin{tcolorbox}[
  ficheSection=colPresence,
  title={Présence dans les sujets analysés (2013-2026)},
]
\begin{center}
\footnotesize\renewcommand{\arraystretch}{1.15}
\begin{tabular}{lccccccccccccc}
\toprule
\textbf{Année} & 13 & 14 & 15 & 16 & 17 & 18 & 19 & 20 & 21 & 22 & 23 & 24 & 25 \\
\midrule
E1 (commune)   & $\checkmark$ & $\checkmark$ & $\checkmark$ & - & $\checkmark$ & $\checkmark$ & - & $\checkmark$ & - & - & $\checkmark$ & $\checkmark$ & $\checkmark$ \\
E2 discipl. IE & - & $\checkmark$ & - & $\checkmark$ & - & - & $\checkmark$ & $\checkmark$ & $\checkmark$ & $\checkmark$ & - & $\checkmark$ & $\checkmark$ \\
\bottomrule
\end{tabular}
\end{center}
\smallskip
\footnotesize\textbf{Sous-thèmes fréquents :}
PD6 - tension moyenne et rapport de transfo (7/13) - facteur de puissance (6/13) - PD2 + hacheur (5/13) - harmoniques côté réseau (4/13) - pont commandé thyristors (2/13)

\medskip
\textbf{Ordres de grandeur (repères rapides).}
\begin{itemize}[leftmargin=1.3em]
  \item Réseau BT FR : tension phase-neutre typique \SI{230}{\volt}, tension phase-phase typique \SI{400}{\volt}, fréquence \SI{50}{\hertz}.
  \item Courants d'étude fréquents en sujet : de quelques \si{\ampere} à quelques dizaines de \si{\ampere}.
  \item En conversion d'énergie, un rendement attendu reste généralement entre 0{,}8 et 0{,}98 selon la technologie.
\end{itemize}

\medskip
\textbf{Synthèse express (à refaire sans support).}
\begin{enumerate}[leftmargin=1.4em]
  \item Poser les hypothèses du modèle et les conventions de signe.
  \item Écrire 2 à 4 relations clés avec unités.
  \item Réaliser un calcul numérique complet et vérifier l'ordre de grandeur.
  \item Justifier le résultat par un argument technique de niveau CAPET.
\end{enumerate}

\end{tcolorbox}

\CAPETENDFILE
